\documentclass[11pt,aspectratio=43]{beamer}
\usepackage{tikz}
\usetheme{Madrid}
\setbeamertemplate{navigation symbols}{}
\title{Constraint propagation: experiments and bounds}
\subtitle{Generated interaction study}
\author{Research group}
\date{}
\year=2026\month=11\day=5
\newcommand{\fuzzlayerA}[1]{\textit{#1}}
\newcommand{\fuzzterm}[1]{\fuzzlayerA{#1}}
\newcommand{\fuzzpair}[2]{\fuzzterm{#1}: #2}
\newcommand{\fuzzunit}{steps}
\let\fuzzoriginalunit\fuzzunit
\newcommand{\fuzzcoefficient}{1}
\newcommand{\fuzzscale}[2][\fuzzcoefficient]{#1\,#2}
\definecolor{fuzzpaint}{RGB}{255,0,255}
\definecolor{fuzzsecond}{RGB}{18,137,67}
\tikzset{fuzzvertex/.style={draw,rounded corners=1pt,fill=blue!10,inner sep=1pt,minimum width=.35cm,minimum height=.25cm}}
\begin{document}\begin{frame}[t,shrink=5]{Pending label}
Before.\par\begin{columns}[onlytextwidth,c]\begin{column}{0.38\textwidth}
\begin{enumerate}
\item<1-> \begin{description}
\item[Stage 1] \begin{tabular}{lrr}Method&Cost&Error\\\hline
Fresh0&24&22\\
Cached1&16&18\\
\end{tabular}
\item[Stage 2] \textrm{A local update} can change the measurements used by several neighboring components.  The estimate is $\frac{a+b}{1+c}$.
\end{description}
\item<2-> \begin{description}
\item[Stage 1] \uncover<2->{\[E=\sum_{i=1}^{n}w_i x_i\]}
\item[Stage 2] \[\begin{pmatrix}a&b\\c&d\end{pmatrix}\begin{pmatrix}x\\y\end{pmatrix}\]
\end{description}
\item<3-> \begin{exampleblock}{Incremental analysis}
\begin{tabular}{lrr}Method&Cost&Error\\\hline
Global0&15&26\\
Local1&10&26\\
Cached2&0&3\\
Fresh3&18&2\\
\end{tabular}
\end{exampleblock}
\end{enumerate}
\end{column}\begin{column}{.58\textwidth}Right text.\end{column}\end{columns}\par Witness.\only<4>{}\end{frame}\end{document}
